\section{Descripción general del proyecto}

El proyecto se desarrolla desde la Ingeniería de Sistemas como una iniciativa académica de apoyo a la gestión del riesgo de incendios forestales en los cerros orientales de Bogotá. Su alcance se concentra en la construcción de un recurso digital orientado a facilitar la interpretación de información sobre el comportamiento del fuego por parte de los organismos de respuesta.

\subsection{Objetivo general}

Desarrollar una aplicación web o móvil para la estimación y visualización del comportamiento probable de incendios forestales en los cerros orientales de Bogotá desde su inicio, a partir de escenarios de simulación precalculados, con el fin de apoyar la toma de decisiones de los bomberos.

\subsection{Objetivos Específicos}

\begin{enumerate}
    \item Construir una biblioteca híbrida de escenarios de incendios forestales propios del territorio de estudio, a partir de simulaciones realizadas con FARSITE y WFDS.
    \item Aprovechar infraestructura HPC e información climática y de viento obtenida mediante una API para la preparación de los escenarios.
    \item Articular los resultados obtenidos con FARSITE y WFDS para representar de manera coherente la propagación y el posible impacto de los incendios forestales en el entorno urbano-forestal de la zona.
    \item Desarrollar una visualización mediante animaciones o videos comprensibles que apoye la interpretación de los escenarios por parte de los bomberos.
\end{enumerate}
